\section{Background}
\label{appx:background}

\subsection{Deferred Definitions}
\label{appx:definitions}
An argument of knowledge allows a computationally bounded prover $\Prover$ to convince a verifier $\Verifier$ that he knows a witness $w$ for a statement~$u$.
\begin{definition}[Argument of Knowledge]
	\label{def: argumentOfKnowledge}
	%The triple
	$(\Setup, \Prover, \Verifier)$ is called an argument of knowledge for the relation $\Relation$ if it satisfies perfect completeness and computational knowledge soundness.
	%as defined below.
\end{definition}

% \cite{groth16_groth2016size}
\begin{definition}[Perfect Completeness]
	\label{def: completeness}
	$(\Setup, \Prover, \Verifier)$ has perfect completeness if for any non-uniform PPT adversary $\Adversary$:
	\begin{equation*}
		% \label{eq:completenessDef}
		\Pr
		\left(
			\left.
				\begin{aligned}
					& (\srs, u, w) \notin \Relation \quad \vee \\
					& \Verifier(\srs, u; \pi) = \TRUE \\
				\end{aligned}
			\right.
			\left |
				\begin{aligned}
					& \srs \leftarrow \Setup(1^{\lambda}); \\
					& (u, w) \leftarrow \Adversary(\srs);\\
					& \pi \leftarrow \PVinteract(\srs, u: w)
				\end{aligned}
			\right.
		\right)
		= 1,
	\end{equation*}
	where $\srs, u, w$, and $\pi$ stand for the structured reference string, the statement, the witness, and the transcript, respectively.
\end{definition}

\begin{definition}[Soundness~\cite{iacr/CanettiCHLRR18}]
	\label{def:soundness}
	$(\Setup, \Prover, \Verifier)$ for a relation $\mathcal{R}$ has a soundness error $\kappa^\mathsf{s}(u)\colon \mathbb{N} \mapsto [0, 1]$ if for any unbounded interactive $\Prover^*$, the probability 
	below
	%that $\Prover^*$ outputs an accepting transcript 
	is bounded by $\kappa^\mathsf{s}(u)$.
	If $\kappa^\mathsf{s}(u)$ is negligible,
	the protocol is sound.
	\begin{equation*}
	% \label{eq:soundnessExpectedTime}
	\Pr
		\left(
			\left.
			\begin{aligned}
				& \Verifier(\pp, u, \pi) = \TRUE \\
			\end{aligned}
			\right.
			\left|
			\begin{aligned}
				& u \notin \Relation;~ \pi \leftarrow (\Prover^* \interact \Verifier) (\pp, u) \\
			\end{aligned}
			\right.
		\right)
	\le \kappa^\mathsf{s}(u).
	\end{equation*}
\end{definition}

\begin{definition}[Knowledge Soundness~\cite{crypto/AttemaCK21}]
	\label{def: soundnessExtract}
	%\quad
	%\phantom{}\\
	$(\Setup, \Prover, \Verifier)$ for relation $\Relation$ is knowledge sound with error $\kappa_{\lambda}(|u|) \colon \mathbb{N} \rightarrow [0, 1]$ if there exists an algorithm $\Extractor$, called a knowledge extractor, given a statement $u$ and black-box oracle access to a (potentially dishonest) prover $\Prover^*$, 
		%the extractor $\Extractor$ 
		it runs in an expected number of steps polynomial in $|u|$ (counting queries to $\Prover$ as a single step) and outputs a witness $ (\srs, u: w) \in \Relation$ with probability: 
	%\begin{equation*}
		% \label{eq:soundnessExpectedTime}
		\[
		\Pr((\srs, u \colon \Extractor(u) ) \in \Relation) 
		\ge
		\frac
		{p(\Prover^*, u) - \kappa_{\lambda}(|u|)}
		{\Poly(|u|)},
		\]
	%\end{equation*}
	where $p(\Prover^*, u)$ is the probability that $\Prover^*$ generates 
	%an accepted
	a valid proof.
\end{definition}



\begin{definition}[Public-Coin]
	\label{def: publicCoin}
	A $(2 \mu + 1)-$move public-coin interactive protocol $(\Setup, \Prover, \Verifier)$ is public-coin if in the $(2j)$-th move, $\Verifier$ sends the $j$-th challenge $x_j \ChooseRandom \mathbb{X}_j$.
	All challenges $x_1, \ldots, x_{\mu}$ sent from the verifier to the prover are public information and chosen uniformly at random and independently of the prover's messages.
\end{definition}

\begin{definition}[Adaptive Fiat-Shamir Transformation]
	The adaptive Fiat-Shamir transformation on a $(2\mu + 1)-$move public-coin interactive protocol: $\mathsf{FS}(\Setup, \Prover, \Verifier) \mapsto (\Setup, \Prover_\mathsf{FS}, \Verifier_{\mathsf{FS}})$ is a two-round argument of knowledge where $\Prover_\mathsf{FS}(\srs, u, w)$ runs $\Prover(\srs, u, w)$ but instead of asking the verifier for the challenge $x_i$ on message $a_i$, the challenges are computed as: 
	\begin{equation*}
		% \label{eq:fiatShamirTransform}
		x_j = \mathsf{Hash}_i(u, a_1, \ldots, a_{j - 1}), \quad j \in [\mu].
	\end{equation*}
the output is then the transcript $\pi = (a_1, \ldots, a_{\mu - 1}, a_{\mu})$.
On input of a statement $u$ and a transcript $\pi = (a_1, \ldots, a_{\mu -1 }, a_{\mu}) $, $\Verifier_{\mathsf{FS}}(\srs, u, \pi)$ accepts if, for $x_j, j \in [\mu]$ as above, $\Verifier$ accepts the transcript $\pi = (a_1, \ldots, a_{\mu - 1}, a_{\mu})$ on input $u$.
\end{definition}


\begin{definition}[Perfect Zero-Knowledge]
\label{def: perfectZeroKnowledge}
	A non-interactive argument of knowledge $(\Setup, \Prover, \Verifier)$ is perfect zero-knowledge~\cite{eurocrypt/Groth16} if there exists a PPT simulator $\Simulator$ and a trapdoor $\mathsf{trap}$ such that for all stateful distinguishers $\Distinguisher$ the following two probabilities are equal:
	\begin{equation*}
		% \label{eq:perfectZeroKnowledge}
		\left.
			\begin{aligned}
			&
			\Pr
			\left(
				\left.
					\begin{aligned}
						& \Distinguisher(\srs, u; \pi; \mathsf{trap}) = \TRUE 
					\end{aligned}
				\right.
				\left |
					\begin{aligned}
						& (\srs, \mathsf{trap}) \leftarrow \Setup(1^{\lambda}); \\
						& \pi \leftarrow \PVinteract (\srs, u: w)
					\end{aligned}
				\right)
			\right.
			\\
			= & 
			\Pr
			\left(
				\left.
					\begin{aligned}
						& \Distinguisher( \srs, u; \pi^*; \mathsf{trap}) = \TRUE \\
					\end{aligned}
				\right.
				\left |
					\begin{aligned}
						& (\srs, \mathsf{trap}) \leftarrow \Setup(1^{\lambda}); \\
						& \pi^* \leftarrow \Simulator (\srs, u; \mathsf{trap})
					\end{aligned}
				\right.
			\right).
			\end{aligned}
		\right.
	\end{equation*}
\end{definition}

\begin{definition}[zkSNARK]
	\label{def:succinctness}
	A zero-knowledge non-interactive argument of knowledge $(\Setup$, $\Prover$, $\Verifier)$ is a succinct non-interactive argument of knowledge (zkSNARK) if:
	\begin{itemize}
	\item The setup time is quasilinear in witness size;
	\item The prover time is quasilinear in witness size;
	\item The proof size is polylogarithmic in witness size;
	\item The verifier time is polylogarithmic in witness size.
	\end{itemize}
\end{definition}



\subsection{Lookup Proof}
\label{appx:lookup}
Lookup proofs allow a prover to show that a collection of queried values,
such as the entries of a matrix $\bm{A}$, are contained in a prescribed
table $\bm{t}\in\Field^m$.
They are a useful primitive in many proof systems:
\eg, they support sparse matrix arithmetization as in
Spartan~\cite{crypto/Setty20}, memory and address consistency in zkVM systems
such as Jolt~\cite{eurocrypt/ArunST24}, and the verification of nonlinear
functions in zero-knowledge machine learning~\cite{ccs/SunLZ24}. Motivated by
these applications, many efficient lookup arguments have been proposed,
including log-derivative lookup (LogUp)~\cite{iacr/Habock22a}, Lasso~\cite{eurocrypt/SettyTW24}, and
Twist and Shout~\cite{iacr/SettyT25}.
We choose LogUp because it scales well for large lookup instances.

\begin{lemma}[Log-Derivative Lookup~{\cite{iacr/Habock22a}}]
\label{lemma:lookup}
Let $\bm{a} \in \Field^N$ and let $\bm{t} \in \Field^m$ be a lookup table. Then $\{a_i\}_{i=1}^N \subseteq \{t_j\}_{j=1}^m$ as multisets iff there exists a multiplicity vector $\bm{\nu} \in \Field^m$ (with $\nu_j$ counting the occurrences of $t_j$ in $\bm{a}$) such that the rational-function identity:
\[
    \sum_{i=1}^N \frac{1}{x + a_i}
    \;=\;
    \sum_{j=1}^m \frac{\nu_j}{x + t_j},
\]
holds in $\Field(x)$. By the Schwartz--Zippel--DeMillo--Lipton lemma, evaluating both sides at a uniformly random $z \sample \Field$ verifies the identity except with probability at most $(N + m)/|\Field|$.
\end{lemma}

\begin{corollary}
Let $\bm{E} \in \Field^{n\times k}$ be a matrix, let
$\bm{t} \in \Field^m$ be a lookup table, and let
$\bm{J} \in [m]^{n\times k}$ be an index matrix.
Let $\bm{m}=(0,1,\ldots,m-1)^\top$.  If $\bm{E}$ is looked up from
$\bm{t}$ at the positions specified by $\bm{J}$, then there exists
$\bm{\nu} \in \Field^m$ such that, for all $\eta,z\in\Field$ for which the
denominators are nonzero:
\[
    \left\langle
    \frac{\bm{1}}{\bm{1}-z(\bm{E}+\eta\bm{J})},
    \bm{1}
    \right\rangle
    =
    \left\langle
    \frac{\bm{1}}{\bm{1}-z(\bm{t}+\eta\bm{m})},
    \bm{\nu}
    \right\rangle,
\]
where division is entrywise. Conversely, if the indexed lookup relation is
false and independent $\eta,z\sample\Field$ are sampled after $\bm E$ is
fixed, the identity holds with probability at most
$O((nk+m)/|\Field|)$.
\end{corollary}

\begin{proof}
The matrix $\bm{E}$ is looked up from $\bm{t}$ at the positions specified by
$\bm{J}$ exactly when every value-index pair $(E_{i,\beta},J_{i,\beta})$ belongs
to the table of pairs $\{(t_j,j)\}_{j=1}^m$, with multiplicities recorded by
$\bm{\nu}$.  The challenge $\eta$ compresses these pairs as:
\[
    E_{i,\beta} + \eta J_{i,\beta}
    \qquad\text{and}\qquad
    t_j + \eta j.
\]
If the pair multisets differ, equality of the compressed multisets is a
nonzero polynomial condition in $\eta$ of degree at most $nk$.  Except with
probability $O(nk/|\Field|)$, the compressed multisets therefore still differ.
Conditioned on this event, Lemma~\ref{lemma:lookup}, evaluated independently at
$z$, accepts with probability at most $(nk+m)/|\Field|$.  The union bound gives
the stated error.
\end{proof}
